\section{Open Challenges and Future Directions}
\label{sec:future}

\revise{The reviewed studies motivate future work on evidence grounding, cross-component reasoning, under-explored subsystems and the capacity to review generated changes.}

\begin{figure}[t!]
\centering
   \includegraphics[width=0.45\textwidth]{figures/challenges.png}
    \caption{Open challenges and future directions for LLMs in OS security.}
    \label{fig:open-challenges}
\end{figure}

\textbf{\revise{Evidence grounding and reproducibility.}}
\revise{LLM outputs are candidates whose security claims need reviewable evidence. Future systems can attach hypotheses to stable traces, code changes and test failures, express rationales as checkable constraints, and report the reproducer, environment, versions and tool settings needed for independent validation. Evaluation should distinguish model suggestions from the checks performed by analyzers, fuzzers or proof assistants and report appropriate baselines and false positives. These artifacts connect model suggestions to the relevant correctness and security properties.}

\textbf{\revise{Cross-component reasoning.}}
\revise{A locally plausible change may violate a contract elsewhere: kernel fixes must preserve concurrency and ABI assumptions, drivers must respect hardware state, and network protocols require stateful oracles. For example, an input can pass through a system call, driver, network stack and file system, crossing ownership and permission boundaries. Retrieval can gather relevant code, specifications and traces without requiring all artifacts in one context; program analysis and specification mining can represent inter-component contracts. Differential tests, sanitizers, model checking and trace queries can then evaluate candidate cross-layer hypotheses.}

\textbf{\revise{Under-explored subsystems.}}
\revise{Hypervisor and TEE interfaces warrant further study because instrumentation and vendor documentation can be limited, while failures may appear only as VM resets or enclave aborts. Promising directions include inferring interface schemas and setup sequences, developing semantic or differential oracles, and evaluating across vendors, versions and configurations. Additional threat models include supply-chain patch risks and side-channel informed inference. Security conclusions should remain specific to the evaluated threat model and evidence.}

\textbf{\revise{Review capacity and vibe coding.}}
\revise{LLM-assisted vibe coding and agentic generate--validate workflows can help maintainers explore candidate fixes, but high submission volumes can place additional demands on review. Proposed directions include summarizing patch intent with supporting traces, identifying concurrency and compatibility risks, and generating targeted regression checks. Clustering related patches and ranking them for human review may support maintainers. Such assistance should be evaluated for review effort, false alarms and missed defects; risks to reviewer capacity are motivations for future evaluation rather than measured outcomes of this survey.}

\revise{The same exploit-development capability can support defensive prioritization and lower the effort required for attacks. Evaluations should therefore measure exploitability and mitigation under explicit threat models while supporting coordinated disclosure and timely review of candidate fixes}~\revcite{openai2026paths,anthropic2026mythos}.
